\documentclass{beamer}

\usepackage[utf8]{inputenc}

\input{../helper}

\title{Introdução à Computação \\ bubble sort}
\author{Alexandre Rademaker}
\date{}

\begin{document}

\begin{frame}
 \maketitle
\end{frame}


\begin{frame}[fragile,allowframebreaks]{Ordenação (sort)}

  Vamos considerar dois algorítmos para esta tarefa:

  \begin{itemize}
  \item selection sort
  \item \emph{bubble sort}
  \item merge sort
  \end{itemize}

\end{frame}


\begin{frame}[fragile,allowframebreaks]{Bubble Sort}

  \begin{lstlisting}[style=code]
    Repeat n - 1 times
      For i from 0 to n - 2
        If numbers[i] and numbers[i + 1] out of order
            Swap them
  \end{lstlisting}

  No \emph{bubble sort} vamos realizar trocas sucessivas de elementos
  adjacentes até que nenhuma troca seja possível.

  Por que o loop mais externo é de $n - 1$ passos?

  \framebreak

  \begin{lstlisting}[style=code]
    5 2 7 4 1 6 3 0
    ^ ^
    2 5 7 4 1 6 3 0
      ^ ^
    2 5 7 4 1 6 3 0
        ^ ^
    2 5 4 7 1 6 3 0
          ^ ^
    2 5 4 1 7 6 3 0
            ^ ^
    2 5 4 1 6 7 3 0
              ^ ^
    2 5 4 1 6 3 7 0
                ^ ^
    2 5 4 1 6 3 0 7
  \end{lstlisting}

  \framebreak

  Agora o maior número está no final da lista, e podemos continuar 

  \begin{lstlisting}[style=code]
   2 5 4 1 6 3 0 | 7
   ^ ^
   2 5 4 1 6 3 0 | 7
     ^ ^
   2 4 5 1 6 3 0 | 7
       ^ ^
   2 4 1 5 6 3 0 | 7
         ^ ^
   2 4 1 5 6 3 0 | 7
           ^ ^
   2 4 1 5 3 6 0 | 7
             ^ ^
   2 4 1 5 3 0 | 6 7
  \end{lstlisting}
  
  A cada passada, a lista de números que precisamos examinar é menor.

  \framebreak

  \begin{lstlisting}[style=normalc,caption=src/bubble.c]
    void bubble_sort (int a[], int n) {
      int i, t, j = n;
      bool s = true;
      
      while (s) {
        s = false;
        for (i = 1; i < j; i++) {
          if (a[i] < a[i - 1]) {
            t = a[i];
            a[i] = a[i - 1];
            a[i - 1] = t;
            s = true;
          }
        }
        j--;
      }
    }
  \end{lstlisting}

  \framebreak

  Epa! Tem alguma coisa diferente em relação ao pseudo-código! O que?

  Ao invés de passar $n - 1$ vezes na lista, passamos enquanto
  trocas forem feitas!

  A troca de valores entre duas posições de um array está ficando
  repetitiva, não?! Vamos falar disso depois\ldots

  \framebreak

  Sobre complexidade? Na versão original

  \begin{align*}
    (n - 1) \times (n - 1) & = \\
    n^2 - 2n + 1 & = O(n^2) 
  \end{align*}
  e também $\Omega(n^2)$.

  Mas com nossa otimização, reduzimos o melhor caso (lista já
  ordenada) para $\Omega(n)$.

\end{frame}


\end{document}

%%% Local Variables:
%%% mode: latex
%%% TeX-master: t
%%% End:
